% 康奈尔听课笔记模板（配色取自王道课件）
\documentclass[10.5pt,a4paper,oneside]{article}
\usepackage[a4paper,top=2.3cm,bottom=2.2cm,left=5.7cm,right=1.5cm,
  marginparwidth=4.25cm,marginparsep=0.55cm,headheight=14pt,footskip=1.1cm]{geometry}
\usepackage{xeCJK}
\usepackage{xeCJKfntef}
\xeCJKDeclareCharClass{FullLeft}{"201C,"2018}
\xeCJKDeclareCharClass{FullRight}{"201D,"2019}
\xeCJKDeclareCharClass{CJK}{"2460->"24FF,"2605,"2606}
\setmainfont{Noto Sans CJK SC}[BoldFont={Noto Sans CJK SC Bold},ItalicFont={Noto Sans CJK SC},ItalicFeatures={FakeSlant=0.18}]
\setsansfont{Noto Sans CJK SC}[BoldFont={Noto Sans CJK SC Bold}]
\setmonofont{DejaVu Sans Mono}[Scale=0.9]
\setCJKmainfont{Noto Sans CJK SC}[BoldFont={Noto Sans CJK SC Bold},ItalicFont={Noto Sans CJK SC},ItalicFeatures={FakeSlant=0.18}]
\setCJKsansfont{Noto Sans CJK SC}[BoldFont={Noto Sans CJK SC Bold}]
\setCJKmonofont{Noto Sans Mono CJK SC}[Scale=0.9]
\linespread{1.32}
\usepackage{amsmath,amssymb}
\usepackage[dvipsnames,table]{xcolor}
\usepackage{graphicx}
\usepackage{adjustbox}
\usepackage{tikz}
\usetikzlibrary{calc}
\usepackage[most]{tcolorbox}
\usepackage{eso-pic}
\usepackage{fancyhdr}
\usepackage{titlesec}
\usepackage{needspace}
\usepackage{enumitem}
\usepackage{xltabular,booktabs,array,colortbl}
\usepackage{marginnote}
\reversemarginpar
\renewcommand*{\raggedleftmarginnote}{}
\renewcommand*{\raggedrightmarginnote}{}
\usepackage{caption}
\usepackage{titletoc}
\usepackage{fancyvrb}
\usepackage{framed}
\usepackage[section]{placeins}
\usepackage{hyperref}
\hypersetup{colorlinks=true,linkcolor=wdblue,urlcolor=wdblue,pdftitle={2.1
课后习题：线性表的定义和基本操作},pdfauthor={听课笔记}}

% ---------- 课件配色 ----------
\definecolor{wdblue}{HTML}{4472C4}
\definecolor{wdlightblue}{HTML}{5B9BD5}
\definecolor{wdtitle}{HTML}{4A86C8}
\definecolor{wddark}{HTML}{1F3864}
\definecolor{wdorange}{HTML}{ED7D31}
\definecolor{wdgray}{HTML}{A5A5A5}
\definecolor{wdgold}{HTML}{FFC000}
\definecolor{wdgreen}{HTML}{70AD47}
\definecolor{wddeepgreen}{HTML}{3F7F2A}
\definecolor{wdred}{HTML}{C00000}
\definecolor{wdpurple}{HTML}{7030A0}
\definecolor{cuebg}{HTML}{F2F6FC}
\definecolor{codebg}{HTML}{F7F9FC}
\color[HTML]{262626}

% ---------- 行内样式 ----------
\newcommand{\hlred}[1]{\textcolor{wdred}{\bfseries #1}}
\newcommand{\hlblue}[1]{\textcolor{wdblue}{\bfseries #1}}
\newcommand{\hlgreen}[1]{\textcolor{wddeepgreen}{\bfseries #1}}
\newcommand{\hlorange}[1]{\textcolor{wdorange}{\bfseries #1}}
\newcommand{\hlpurple}[1]{\textcolor{wdpurple}{\bfseries #1}}
\newcommand{\uln}[1]{\CJKunderline[format=\color{wdred},thickness=0.9pt,depth=0.28em]{#1}}
\newcommand{\wavy}[1]{\CJKunderwave[format=\color{wdorange}]{#1}}
\newcommand{\hlmark}[1]{\CJKunderline[format=\color{wdgold!55},thickness=0.55em,depth=-0.05em]{#1}}
\newcommand{\tagbox}[2]{\tcbox[on line,colback=#1,colframe=#1,coltext=white,boxsep=0.5pt,left=3pt,right=3pt,top=0.5pt,bottom=0.5pt,arc=2pt,fontupper=\small\bfseries]{#2}}

% ---------- 康奈尔版式背景：左侧线索栏 ----------
\newlength{\cuestripw}\setlength{\cuestripw}{5.15cm}
\AddToShipoutPictureBG{%
  \begin{tikzpicture}[remember picture,overlay]
    \fill[cuebg] (current page.north west) rectangle ($(current page.south west)+(\cuestripw,0)$);
    \draw[wdblue,line width=1.2pt] ($(current page.north west)+(\cuestripw,-1.55cm)$) -- ($(current page.south west)+(\cuestripw,1.55cm)$);
    \node[anchor=west,text=wdblue,font=\small\bfseries] at ($(current page.north west)+(0.75cm,-1.25cm)$) {线索 · 问题};
    \node[anchor=west,text=wdblue,font=\small\bfseries] at ($(current page.north west)+(\cuestripw+0.55cm,-1.25cm)$) {笔记 · 老师主线 + 补充讲解};
  \end{tikzpicture}}

% ---------- 页眉页脚 ----------
\pagestyle{fancy}\fancyhf{}
\renewcommand{\headrulewidth}{0pt}
\fancyhead[R]{\small\color{wdgray}第2章 线性表}
\fancyfoot[C]{\small\color{wdgray}\thepage}
\fancyfoot[R]{\small\color{wdorange}2.1
课后习题：线性表的定义和基本操作 · 听课笔记}

\fancypagestyle{plain}{\fancyhf{}\fancyfoot[C]{\small\color{wdgray}\thepage}\fancyfoot[R]{\small\color{wdorange}2.1
课后习题：线性表的定义和基本操作 · 听课笔记}}

% ---------- 标题：仿课件的蓝色标题 + 六色条 ----------
\newcommand{\sixbar}{\par\nointerlineskip\vspace{3pt}\centering
  \begin{tikzpicture}\foreach \c [count=\i] in {wdlightblue,wdorange,wdgray,wdgold,wdblue,wdgreen}
     {\fill[\c] ({(\i-1)*1.05},0) rectangle ({(\i-1)*1.05+0.92},0.11);}\end{tikzpicture}\par}
\titleformat{\section}[block]{\centering\Large\bfseries\color{wdtitle}}{}{0pt}{}[\sixbar]
\titlespacing*{\section}{0pt}{16pt}{10pt}
\pretocmd{\section}{\Needspace{8\baselineskip}}{}{}
\titleformat{\subsection}{\Needspace{7\baselineskip}\large\bfseries\color{wddark}}{\textcolor{wdblue}{\rule[-0.15em]{4pt}{1.05em}}\hspace{6pt}}{0pt}{}
\titlespacing*{\subsection}{0pt}{12pt}{5pt}
\titleformat{\subsubsection}{\Needspace{5\baselineskip}\normalsize\bfseries\color{wdblue}}{\textcolor{wdorange}{\large\ensuremath{\blacktriangleright}}\hspace{4pt}}{0pt}{}
\titlespacing*{\subsubsection}{0pt}{8pt}{3pt}
\setcounter{secnumdepth}{0}
\setlength{\parindent}{0pt}\setlength{\parskip}{5pt}
\setlist{nosep,leftmargin=1.5em,itemsep=2pt}
\setlist[itemize,1]{label=\textcolor{wdblue}{\small\ensuremath{\bullet}}}
\setlist[itemize,2]{label=\textcolor{wdorange}{\scriptsize\ensuremath{\blacktriangleright}}}
\setlist[enumerate,1]{label=\textcolor{wdblue}{\bfseries\arabic*.}}
\providecommand{\tightlist}{\setlength{\itemsep}{2pt}\setlength{\parskip}{0pt}}

% ---------- 目录 ----------
\setcounter{tocdepth}{3}
\titlecontents{section}[0pt]{\addvspace{5pt}\bfseries\color{wddark}}{}{}{\titlerule*[0.7pc]{\color{wdgray!60}.}\contentspage}
\titlecontents{subsection}[1.6em]{\addvspace{1pt}\small\color{black!75}}{}{}{\titlerule*[0.7pc]{\color{wdgray!40}.}\contentspage}
\newcommand{\tocback}{\par\vspace{-3pt}{\hfill\footnotesize\hyperlink{toc}{\textcolor{wdgray}{\ensuremath{\uparrow}\,回到目录}}}\par\vspace{-4pt}}

% ---------- 线索栏 ----------
\newcommand{\cueblock}[2]{\Needspace{4\baselineskip}\leavevmode#1{#2}\par\vspace{-\baselineskip}\vspace{-\parskip}}
% 线索直接贴在段首行的左侧（不用 marginnote，避免分页时偶尔落进正文栏）
\newcommand{\cuebox}[1]{\leavevmode\llap{\smash{\begin{minipage}[t]{\marginparwidth}\raggedright\small\linespread{1.18}\selectfont #1\end{minipage}}\hspace{\marginparsep}}}
\newcommand{\cuenote}[1]{\cuebox{{\color{wdblue}\bfseries Q\,}\color{wddark}#1}}
\newcommand{\cuekey}[1]{\cuebox{\color{wdred}\bfseries #1}}

% ---------- 讲解框 ----------
\tcbset{wdbox/.style={enhanced,breakable,arc=2pt,boxrule=0pt,leftrule=3pt,left=7pt,right=7pt,top=5pt,bottom=5pt,
   before skip=7pt,after skip=7pt,fonttitle=\bfseries\small,attach title to upper={\quad},coltitle=#1,
   colframe=#1,colback=#1!6}}
\newtcolorbox{defbox}[1][]{enhanced,breakable,arc=3pt,boxrule=0.8pt,colframe=wdblue,colback=wdblue!4,
   left=8pt,right=8pt,top=6pt,bottom=6pt,fonttitle=\bfseries\small,coltitle=white,colbacktitle=wdblue,
   title={定义#1},before skip=8pt,after skip=8pt}
\newtcolorbox{bridgebox}{wdbox=wdblue,title={承上启下}}
\newtcolorbox{plainbox}{wdbox=wddeepgreen,colframe=wdgreen,colback=wdgreen!8,title={大白话}}
\newtcolorbox{deepbox}{wdbox=wdpurple,title={背后的思想}}
\newtcolorbox{exambox}{wdbox=wdred,title={考点}}
\newtcolorbox{pitfallbox}{wdbox=wdorange,title={易错提醒}}
\newtcolorbox{bookbox}{wdbox=wddark,colback=wdgray!10,colframe=wdgray,title={王道书补充}}
\newtcolorbox{teacherbox}{wdbox=wdorange,colback=wdgold!8,colframe=wdgold,coltitle=wdorange,title={老师原话}}
\newtcolorbox{fixbox}{wdbox=wdred,colback=wdred!4,title={勘误 / 说明}}
% 习题笔记用
\newtcolorbox{questionbox}[1][]{enhanced,breakable,arc=3pt,boxrule=0.8pt,colframe=wddark,colback=white,
   left=8pt,right=8pt,top=6pt,bottom=6pt,fonttitle=\bfseries\small,coltitle=white,colbacktitle=wddark,
   title={#1},before skip=10pt,after skip=6pt}
\newtcolorbox{videobox}{wdbox=wdorange,colback=wdgold!8,colframe=wdgold,coltitle=wdorange,title={视频讲解}}
\newtcolorbox{jiexibox}{wdbox=wddark,colback=wdgray!10,colframe=wdgray,title={书上解析}}
\newtcolorbox{addbox}{wdbox=wdpurple,title={我的补充}}
\newtcolorbox{recallbox}{wdbox=wdblue,title={回顾正课}}
\newtcolorbox{watchbox}{enhanced,breakable,arc=3pt,boxrule=1pt,colframe=wdorange,colback=wdgold!12,
   fonttitle=\bfseries\small,coltitle=white,colbacktitle=wdorange,title={建议去看视频},left=8pt,right=8pt,before skip=8pt,after skip=8pt}
\newtcolorbox{summarybox}{enhanced,arc=3pt,
   boxrule=1pt,colframe=wdblue,colback=white,colbacktitle=wdblue,coltitle=white,fonttitle=\bfseries,
   title={总结区 · 用自己的话复述这一节},left=10pt,right=10pt,top=8pt,bottom=8pt,before skip=14pt}

% ---------- 代码 ----------

\makeatletter
\@ifundefined{Shaded}{}{\renewenvironment{Shaded}{\def\FrameCommand{\fboxsep=5pt\fboxrule=0.5pt\fcolorbox{wdblue!35}{codebg}}%
   \MakeFramed{\advance\hsize-\width\FrameRestore}\small}{\endMakeFramed}}
\makeatother
\fvset{fontsize=\small}

% ---------- 表格 ----------
\renewcommand{\arraystretch}{1.42}
\newcolumntype{Y}{>{\raggedright\arraybackslash}X}

% ---------- 图片 ----------
% 课件截图做成浮动体：LaTeX 会自动用正文填满页面，避免大块空白
\setcounter{topnumber}{3}\setcounter{bottomnumber}{2}\setcounter{totalnumber}{4}
\renewcommand{\topfraction}{0.85}\renewcommand{\bottomfraction}{0.6}\renewcommand{\textfraction}{0.12}\renewcommand{\floatpagefraction}{0.75}
\setlength{\textfloatsep}{10pt plus 3pt minus 2pt}\setlength{\floatsep}{8pt plus 2pt}
\newcommand{\slidefig}[3]{\begin{figure}[!ht]\centering\begin{tcolorbox}[enhanced,width=#2\linewidth,arc=2pt,boxrule=0.6pt,colframe=wdgray!70,
   colback=white,left=2pt,right=2pt,top=2pt,bottom=2pt,drop fuzzy shadow=black!25]\includegraphics[width=\linewidth]{#1}\end{tcolorbox}
   \vspace{-3pt}{\small\color{wdgray}#3}\end{figure}}
\newcommand{\plainfig}[3]{\par\noindent\begin{minipage}{\linewidth}\centering\includegraphics[width=#2\linewidth]{#1}\par\vspace{2pt}{\small\color{wdgray}#3}\end{minipage}\par\vspace{4pt}}

\begin{document}
% ---------- 封面条（仿课件封面） ----------
\thispagestyle{plain}
\begin{tcolorbox}[enhanced,enlarge left by=-4.8cm,width=\linewidth+4.8cm,arc=0pt,boxrule=0pt,
   colback=wddark,interior style={left color=wddark,right color=wdblue},left=18pt,right=18pt,top=14pt,bottom=12pt,
   overlay={\draw[white,line width=2.5pt] ([xshift=10pt,yshift=-10pt]frame.north west) -- ++(0,-0.9cm);
            \draw[white,line width=2.5pt] ([xshift=10pt,yshift=-10pt]frame.north west) -- ++(1.6cm,0);
            \draw[white,line width=2.5pt] ([xshift=-10pt,yshift=10pt]frame.south east) -- ++(0,0.9cm);
            \draw[white,line width=2.5pt] ([xshift=-10pt,yshift=10pt]frame.south east) -- ++(-1.6cm,0);}]
  {\color{white!85!wdlightblue}\small 数据结构 · 第2章
线性表 · 康奈尔听课笔记}\par\vspace{6pt}
  {\color{white}\fontsize{22}{28}\selectfont\bfseries 2.1
课后习题：线性表的定义和基本操作}\par\vspace{6pt}
  {\color{white!80!wdlightblue}\small 王道书 §2.1.3 试题精选 · §2.1.4
答案与解析（第 14--15 页）· 习题讲解字幕 P6（选择 1--4）}
\end{tcolorbox}
\vspace{4pt}
\hypertarget{ux8fd9ux7ec4ux9898ux5728ux8003ux4ec0ux4e48}{%
\section{这组题在考什么}\label{ux8fd9ux7ec4ux9898ux5728ux8003ux4ec0ux4e48}}

\Needspace{4\baselineskip}

\cuenote{4 道题都在考哪一句话？}4 道题全是定义题，考的就是正课 2.1
那句定义和前驱后继的说法：

\begin{defbox}

线性表是具有\hlred{相同数据类型}的 \(n\,(n\ge0)\)
个\hlred{数据元素}的\hlred{有限}\hlred{序列}。除第一个元素外，每个元素有且仅有一个直接前驱；除最后一个元素外，每个元素有且仅有一个直接后继。

\end{defbox}

{\rowcolors{2}{wdblue!5}{white}\arrayrulecolor{wdblue!35}
\begin{xltabular}{\linewidth}{>{\hsize=1.000\hsize\raggedright\arraybackslash}X>{\hsize=1.000\hsize\raggedright\arraybackslash}X>{\hsize=1.000\hsize\raggedright\arraybackslash}X}
\rowcolor{wdblue}\textcolor{white}{\bfseries 题号} & \textcolor{white}{\bfseries 考点} & \textcolor{white}{\bfseries 答案} \\
\endfirsthead
01 & 线性表由什么组成 & C \\ \hline
02 & 哪个描述是线性表 & B \\ \hline
03 & 前驱、后继的个数 & A \\ \hline
04 & 没有前驱也没有后继的情况 & A \\ \hline
\end{xltabular}\par\vspace{4pt}}

\begin{fixbox}

\textbf{这节的习题讲解视频不用看。} 4
道题都只是复述定义，没有正课和书上之外的新内容，也没有需要截图的图。视频里有两处口误，下面对应的题里已经标出。

\end{fixbox}

\Needspace{8\baselineskip}\vspace{22pt}

\cueblock{\cuenote}{\relax{线性表的元素是数据元素还是数据项？}}

\begin{questionbox}[01]

线性表是具有 \(n\) 个（　）的有限序列。

A. 数据表　　　　B. 字符　　　　C. 数据元素　　　　D. 数据项

\end{questionbox}

\hlred{\textbf{答案：C}}

\begin{videobox}

直接套定义。A 数据表只是把元素排成表格的一种组织形式；B
字符只是一种数据类型，“\(n\)
个字符的有限序列是线性表”成立，反过来“线性表是 \(n\)
个字符的有限序列”就不成立了；D
数据项只是数据元素的组成部分，比如学生的学号、姓名、性别。

\end{videobox}

\begin{jiexibox}

线性表是由具有相同数据类型的有限数据元素组成的，数据元素是由数据项组成的。

\end{jiexibox}

\begin{fixbox}

习题讲解视频里说数据项是“构成数据元素的\textbf{基本单位}”，这是口误。按正课
1.1：\hlred{数据元素是基本单位，数据项是最小单位}（不可再分）。

\end{fixbox}

\Needspace{8\baselineskip}\vspace{22pt}

\cueblock{\cuenote}{\relax{集合、无限序列、存储结构，为什么都不是线性表？}}

\begin{questionbox}[02]

下列几种描述中，（　）是一个线性表。

A. 由 \(n\) 个实数组成的集合\\
B. 由 100 个字符组成的序列\\
C. 所有整数组成的序列\\
D. 邻接表

\end{questionbox}

\hlred{\textbf{答案：B}}

\begin{videobox}

抓住“相同数据类型”和“有限序列”两个关键点，逐个排除：A
同类型、有限，但它是\textbf{集合}，集合里的元素没有先后，没有前驱后继；C
所有整数是无限的；D 邻接表是图的一种存储结构，根本不是逻辑结构。只有 B
三条全满足。

\end{videobox}

\begin{jiexibox}

线性表定义的要求为：相同数据类型、有限序列。C
的元素个数是无穷个，错误；A 集合中的元素没有前后驱关系，错误；D
属于一种存储结构，本题要求选出的是一个具体的线性表，不要将二者混为一谈。

\end{jiexibox}

\begin{recallbox}

三个错误选项正好各缺定义里的一条：

{\rowcolors{2}{wdblue!5}{white}\arrayrulecolor{wdblue!35}
\begin{xltabular}{\linewidth}{>{\hsize=1.143\hsize\raggedright\arraybackslash}X>{\hsize=0.952\hsize\raggedright\arraybackslash}X>{\hsize=0.952\hsize\raggedright\arraybackslash}X>{\hsize=0.952\hsize\raggedright\arraybackslash}X}
\rowcolor{wdblue}\textcolor{white}{\bfseries 选项} & \textcolor{white}{\bfseries 同类型} & \textcolor{white}{\bfseries 有限} & \textcolor{white}{\bfseries 序列（有先后）} \\
\endfirsthead
A. \(n\) 个实数的集合 & ✓ & ✓ & \hlred{✗} \\ \hline
B. 100 个字符的序列 & ✓ & ✓ & ✓ \\ \hline
C. 所有整数的序列 & ✓ & \hlred{✗} & ✓ \\ \hline
\end{xltabular}\par\vspace{4pt}}

D 是另一层面的错误。正课 2.1
笔记里书上的注意：线性表是\textbf{逻辑结构}，顺序表、链表是\textbf{存储结构}；邻接表和它们一样是存储结构（图的，第
6 章学）。C 选项就是正课里老师举的那个反例。

\end{recallbox}

\Needspace{8\baselineskip}\vspace{22pt}

\cueblock{\cuenote}{\relax{除开始元素外，每个元素有几个前驱？}}

\begin{questionbox}[03]

在线性表中，除开始元素外，每个元素（　）。

A. 只有唯一的前驱元素　　　　B. 只有唯一的后继元素\\
C. 有多个前驱元素　　　　　　D. 有多个后继元素

\end{questionbox}

\hlred{\textbf{答案：A}}

\begin{videobox}

举例 \((A,B,C,D)\)：除了开头的 \(A\)，\(B\)、\(C\)、\(D\)
都有且仅有一个前驱。B 不对，因为最后一个元素 \(D\) 没有后继；C、D
不对，线性表是一对一的关系，不存在多个前驱或多个后继。

\end{videobox}

\begin{jiexibox}

线性表中，除最后一个（或第一个）元素外，每个元素都只有一个后继（或前驱）元素。

\end{jiexibox}

\begin{fixbox}

习题讲解视频里老师复述书上定义时说成“除了第一个元素之外，每个元素有且仅有一个直接\textbf{后继}”，这是口误：除第一个元素外的是\uln{直接前驱}，除最后一个元素外的才是直接后继。他后面举例时用的是对的。

\end{fixbox}

\begin{addbox}

做这类题先对“开始”还是“最后”：\textbf{开始元素没有前驱，最后元素没有后继}。题干说“除开始元素外”，就是在问前驱。

\end{addbox}

\Needspace{8\baselineskip}\vspace{22pt}

\cueblock{\cuenote}{\relax{既没有前驱也没有后继，说明表里有几个元素？}}

\begin{questionbox}[04]

若非空线性表中的元素既没有直接前驱，又没有直接后继，则该表中有（　）个元素。

A. 1　　　　B. 2　　　　C. 3　　　　D. \(n\)

\end{questionbox}

\hlred{\textbf{答案：A}}

\begin{videobox}

只要有两个及以上的元素，第一个就有后继，第二个就有前驱。既没前驱又没后继，只能是孤零零的一个元素。

\end{videobox}

\begin{jiexibox}

线性表中的第一个元素没有直接前驱，最后一个元素没有直接后继；当线性表中仅有一个元素时，该元素既没有直接前驱，又没有直接后继。

\end{jiexibox}

\begin{addbox}

仅有的这个元素既是表头元素 \(a_1\)，也是表尾元素 \(a_n\)（此时
\(n=1\)）。题干的“非空”是为了排除空表，空表里根本没有元素可谈。

\end{addbox}

\begin{summarybox}

\textbf{这 4 题考的都是定义。}
线性表的元素是数据元素（不是数据项、不是某种具体类型），要同类型、有限、有先后；集合缺“先后”，无限序列缺“有限”，邻接表、顺序表、链表是存储结构，和线性表这个逻辑结构不在一个层面。前驱后继：开头没前驱、末尾没后继，其余各有且仅有一个；只有一个元素时两样都没有。

\textbf{自测}：① “由 \(n\) 个整数组成的集合”为什么不是线性表？② 线性表
\((a_1)\) 中 \(a_1\) 是表头还是表尾？

\end{summarybox}
\end{document}
